% ============================================================
% Proposition 3.4: Single-Domain Equilibria
% Status: FULL RECONSTRUCTION
% ============================================================

% --- Definitions and Assumptions ---
\begin{definition}[Single-Domain State Space]
When $|D| = 1$, denote $\sigma_A = \sigma_A(d, \cdot)$,
$\delta_A = \delta_A(d, \cdot)$, $\alpha_A = \alpha_A(d, \cdot)$.
The reduced ODE system is:
\begin{align*}
\dot{\delta}_A &= f_{\text{learn}} \eta_{\max} e^{-a e^{-b \delta_A}} T_A
- \lambda_c (1 - \gamma_\sigma \sigma_A) \delta_A (1 - r_A), \\
\dot{\sigma}_A &= \sigma_A [\rho P_A \alpha_A (1 - \gamma_\sigma \sigma_A) - \epsilon_\sigma \Omega_{SL}], \\
\dot{\alpha}_A &= \gamma C_A (1 - \alpha_A) - \zeta_\alpha \alpha_A R_A^{\text{surface}}.
\end{align*}
\end{definition}

\begin{assumption}[Symmetric Couplings]
All cross-domain coupling terms vanish when $|D| = 1$: no intersection
discovery, no interdomain schema transfer, no multi-domain attention.
\end{assumption}

\begin{assumption}[Non-Degenerate Parameters]
The parameters satisfy:
\begin{align*}
f_{\text{learn}} &> 0, \quad \eta_{\max} > 0, \quad \lambda_c > 0, \\
\rho &> 0, \quad P_A > 0, \quad \gamma_\sigma > 0, \quad \epsilon_\sigma > 0, \quad \Omega_{SL} > 0, \\
\gamma &> 0, \quad C_A > 0, \quad \zeta_\alpha > 0, \quad R_A^{\text{surface}} > 0.
\end{align*}
\end{assumption}

% --- Proposition Statement ---
\begin{proposition}[Single-Domain Equilibria]
\label{prop:equilibria}
Under the single-domain reduction (Assumptions~1--2), the system has
precisely two equilibrium families:
\begin{enumerate}[nosep]
\item \textbf{Schema-free equilibrium $\mathcal{E}_0$:}
  $\sigma_A = 0$, $\delta_A = 0$,
  $\alpha_A^* = \dfrac{\gamma C_A}{\gamma C_A + \zeta_\alpha R_A^{\text{surface}}}$.
  Exists for all parameter values.

\item \textbf{Schema-active equilibrium $\mathcal{E}_1$:}
  \begin{align*}
  \sigma_A^* &= \frac{1}{\gamma_\sigma}\left(1 - \frac{1}{R_0}\right), \\
  \alpha_A^* &= \frac{\gamma C_A}{\gamma C_A + \zeta_\alpha R_A^{\text{surface}}},\\
  \delta_A^* &: f_{\text{learn}} \eta(\delta_A^*) T_A = \lambda_c (1 - \gamma_\sigma \sigma_A^*) \delta_A^* (1 - r_A),
  \end{align*}
  where $R_0 = \rho P_A \alpha_A^* / \epsilon_\sigma \Omega_{SL}$.
  Exists iff $R_0 \geq 1$.
\end{enumerate}
The stability of each equilibrium is determined by the eigenvalues of
the $3\times 3$ Jacobian, with $\mathcal{E}_0$ stable for $R_0 < 1$
and $\mathcal{E}_1$ stable for $R_0 > 1$.
\end{proposition}

% --- Proof ---
\begin{proof}

\paragraph{Step 1: Equilibrium equations.}
Setting $\dot{\delta}_A = \dot{\sigma}_A = \dot{\alpha}_A = 0$:

For $\dot{\sigma}_A = 0$, the product form gives either:
\[
\sigma_A = 0 \quad\text{or}\quad \rho P_A \alpha_A (1 - \gamma_\sigma \sigma_A) - \epsilon_\sigma \Omega_{SL} = 0.
\]

For $\dot{\alpha}_A = 0$:
\[
\gamma C_A (1 - \alpha_A) = \zeta_\alpha \alpha_A R_A^{\text{surface}}
\;\Longrightarrow\; \alpha_A^* = \frac{\gamma C_A}{\gamma C_A + \zeta_\alpha R_A^{\text{surface}}},
\]
which is independent of $\sigma_A$ and $\delta_A$. For positive
parameters, $\alpha_A^* \in (0,1)$, confirming the attractiveness
of the attentional subsystem to a unique fixed point.

For $\dot{\delta}_A = 0$:
\[
f_{\text{learn}} \eta(\delta_A) T_A = \lambda_c (1 - \gamma_\sigma \sigma_A) \delta_A (1 - r_A).
\]
This equation defines $\delta_A^*$ implicitly; the left-hand side is
monotonically increasing in $\delta_A$ (since $\eta' > 0$) and the
right-hand side is linear in $\delta_A$. A positive solution exists
iff the left-hand side exceeds the right-hand side at $\delta_A = 0^+$,
which holds for all positive parameters.

\paragraph{Step 2: Schema-free equilibrium $\mathcal{E}_0$.}
At $\sigma_A = 0$, the depth equilibrium condition becomes:
\[
f_{\text{learn}} \eta(\delta_A) T_A = \lambda_c \delta_A (1 - r_A),
\]
which has the trivial solution $\delta_A = 0$ (no depth accumulation
in the absence of schema coherence). The full equilibrium is:
\[
\mathcal{E}_0 = (\delta_A = 0,\; \sigma_A = 0,\; \alpha_A = \alpha_A^*).
\]

\paragraph{Step 3: Schema-active equilibrium $\mathcal{E}_1$.}
Substituting $\alpha_A^*$ into the non-trivial $\dot{\sigma}_A = 0$
condition:
\[
\rho P_A \alpha_A^* (1 - \gamma_\sigma \sigma_A) = \epsilon_\sigma \Omega_{SL}.
\]
Define the basic reproduction number:
\[
R_0 = \frac{\rho P_A \alpha_A^*}{\epsilon_\sigma \Omega_{SL}}.
\]
Then:
\[
1 - \gamma_\sigma \sigma_A^* = \frac{1}{R_0}
\;\Longrightarrow\; \sigma_A^* = \frac{1}{\gamma_\sigma}\left(1 - \frac{1}{R_0}\right).
\]
For $\sigma_A^* > 0$, we require $R_0 > 1$. At $R_0 = 1$, the
schema-active equilibrium coincides with the schema-free equilibrium
(transcritical bifurcation; see Proposition~4.1).

The depth $\delta_A^*$ is the positive solution of:
\[
f_{\text{learn}} \eta(\delta_A^*) T_A = \lambda_c (1 - \gamma_\sigma \sigma_A^*) \delta_A^* (1 - r_A)
= \frac{\lambda_c}{R_0} \delta_A^* (1 - r_A),
\]
where the decay coefficient $\lambda_c / R_0$ is smaller than
$\lambda_c$, confirming that schema-active agents retain depth
more effectively.

\paragraph{Step 4: $3\times 3$ Jacobian at equilibrium.}
The Jacobian $\mathbf{J} \in \mathbb{R}^{3\times 3}$ has entries
$J_{ij} = \partial \dot{x}_i / \partial x_j$ for
$\mathbf{x} = (\delta_A, \sigma_A, \alpha_A)$:

\[
\mathbf{J} =
\begin{pmatrix}
J_{\delta\delta} & J_{\delta\sigma} & 0 \\
0 & J_{\sigma\sigma} & J_{\sigma\alpha} \\
0 & 0 & J_{\alpha\alpha}
\end{pmatrix},
\]
where:
\begin{align*}
J_{\delta\delta} &= \frac{\partial \dot{\delta}_A}{\partial \delta_A}
= \frac{f_{\text{learn}} T_A \cdot \eta_{\max} a b e^{-b\delta_A} e^{-a e^{-b\delta_A}}}
{- \lambda_c (1 - \gamma_\sigma \sigma_A)(1 - r_A)}, \\
J_{\delta\sigma} &= \frac{\partial \dot{\delta}_A}{\partial \sigma_A}
= \lambda_c \gamma_\sigma \delta_A (1 - r_A), \\
J_{\sigma\sigma} &= \frac{\partial \dot{\sigma}_A}{\partial \sigma_A}
= \rho P_A \alpha_A (1 - 2\gamma_\sigma \sigma_A) - \epsilon_\sigma \Omega_{SL}, \\
J_{\sigma\alpha} &= \frac{\partial \dot{\sigma}_A}{\partial \alpha_A}
= \rho P_A \sigma_A (1 - \gamma_\sigma \sigma_A), \\
J_{\alpha\alpha} &= \frac{\partial \dot{\alpha}_A}{\partial \alpha_A}
= -\gamma C_A - \zeta_\alpha R_A^{\text{surface}} < 0.
\end{align*}

The lower-triangular block structure is due to the fact that
$\dot{\alpha}_A$ does not depend on $\delta_A$ or $\sigma_A$, and
$\dot{\sigma}_A$ does not depend on $\delta_A$.

\paragraph{Step 5: Stability of $\mathcal{E}_0$ ($\sigma_A = 0$).}
At $\mathcal{E}_0$, $\sigma_A = 0$, $\delta_A = 0$:
\begin{align*}
J_{\delta\delta} &= f_{\text{learn}} T_A \eta_{\max} e^{-a} - \lambda_c (1 - r_A), \\
J_{\delta\sigma} &= 0 \quad (\text{since } \delta_A = 0), \\
J_{\sigma\sigma} &= \rho P_A \alpha_A^* - \epsilon_\sigma \Omega_{SL}
= \epsilon_\sigma \Omega_{SL} (R_0 - 1), \\
J_{\sigma\alpha} &= 0 \quad (\text{since } \sigma_A = 0), \\
J_{\alpha\alpha} &= -\gamma C_A - \zeta_\alpha R_A^{\text{surface}} < 0.
\end{align*}

The eigenvalues are the diagonal entries:
\[
\lambda_1 = J_{\delta\delta},\quad
\lambda_2 = J_{\sigma\sigma} = \epsilon_\sigma \Omega_{SL} (R_0 - 1),\quad
\lambda_3 = J_{\alpha\alpha} < 0.
\]

At $\delta_A = 0$, the learning function satisfies
$\eta(0) = \eta_{\max} e^{-a}$, so:
\[
\lambda_1 = f_{\text{learn}} T_A \eta_{\max} e^{-a} - \lambda_c (1 - r_A).
\]
For small initial depth accumulation, the learning drive exceeds the
decay ($\lambda_1 > 0$), so the $\delta_A$ direction is unstable at
$\delta_A = 0$ --- depth grows away from zero, consistent with
Proposition~3.1's local existence result that $\delta_A$ increases
from initialisation.

The critical eigenvalue is $\lambda_2$:
\begin{itemize}
\item If $R_0 < 1$: $\lambda_2 < 0$, so $\mathcal{E}_0$ is stable in
  the $\sigma_A$ direction. Schema perturbations decay.
\item If $R_0 > 1$: $\lambda_2 > 0$, so $\mathcal{E}_0$ is unstable
  in the $\sigma_A$ direction. Schema perturbations grow, driving the
  system toward $\mathcal{E}_1$.
\item If $R_0 = 1$: $\lambda_2 = 0$, the non-hyperbolic transcritical
  bifurcation point (Proposition~4.1).
\end{itemize}

Thus $\mathcal{E}_0$ is a saddle for $R_0 > 1$ and an attractor
(in the $\sigma_A$-$\alpha_A$ subspace) for $R_0 < 1$.

\paragraph{Step 6: Stability of $\mathcal{E}_1$ ($\sigma_A^* > 0$).}
At $\mathcal{E}_1$, $\sigma_A^*$ satisfies $1 - \gamma_\sigma \sigma_A^* = 1/R_0$.
Substituting into $J_{\sigma\sigma}$:
\[
J_{\sigma\sigma} = \rho P_A \alpha_A^* (1 - 2\gamma_\sigma \sigma_A^*) - \epsilon_\sigma \Omega_{SL}
= \rho P_A \alpha_A^* \left(1 - 2\left(1 - \frac{1}{R_0}\right)\right) - \epsilon_\sigma \Omega_{SL}
= \rho P_A \alpha_A^* \left(\frac{2}{R_0} - 1\right) - \epsilon_\sigma \Omega_{SL}.
\]

Using $R_0 = \rho P_A \alpha_A^* / \epsilon_\sigma \Omega_{SL}$:
\[
J_{\sigma\sigma} = \epsilon_\sigma \Omega_{SL} \left[R_0\left(\frac{2}{R_0} - 1\right) - 1\right]
= \epsilon_\sigma \Omega_{SL} (1 - R_0).
\]

For $R_0 > 1$, $J_{\sigma\sigma} < 0$: the $\sigma_A$ direction is
stable at $\mathcal{E}_1$.

The $\alpha_A$ eigenvalue $J_{\alpha\alpha} < 0$ is unchanged.

The $\delta_A$ eigenvalue $J_{\delta\delta}$ at $\mathcal{E}_1$ is:
\[
J_{\delta\delta} = f_{\text{learn}} T_A \eta'(\delta_A^*) - \frac{\lambda_c}{R_0} (1 - r_A),
\]
where $\eta'(\delta_A^*) = \eta_{\max} a b e^{-b\delta_A^*} e^{-a e^{-b\delta_A^*}}$.
Since $\delta_A^*$ is the fixed point of the balanced equation,
$J_{\delta\delta} < 0$ by the monotonicity of the learning saturation
function (the right-hand side of the $\delta_A$ ODE crosses zero from
above at $\delta_A^*$, making the equilibrium stable in the $\delta_A$
direction).

Therefore, for $R_0 > 1$, all three eigenvalues at $\mathcal{E}_1$
have negative real parts: the schema-active equilibrium is locally
asymptotically stable.
$\square$
\end{proof}

% --- Boundary Cases ---
\begin{boundary-case}
\textbf{$R_0 = 1$ (bifurcation threshold):}
Both equilibria coincide at $\sigma_A = 0$. The Jacobian has a zero
eigenvalue, and the system undergoes a transcritical bifurcation
(formalised in Proposition~4.1). Higher-order terms determine the
stability beyond the linear approximation.

\textbf{$\gamma_\sigma = 0$ (no schema modulation of depth decay):}
The depth decay coefficient does not depend on $\sigma_A$, so
$J_{\delta\sigma} = 0$ at both equilibria. The $\delta_A$ and
$\sigma_A$ subsystems decouple linearly. Equilibrium existence
conditions remain unchanged.

\textbf{$\alpha_A^* \to 1$ (maximal attention):}
Occurs when $\gamma C_A \gg \zeta_\alpha R_A^{\text{surface}}$.
The reproduction number $R_0$ approaches its maximum
$\rho P_A / \epsilon_\sigma \Omega_{SL}$, maximising the basin
of the schema-active equilibrium.

\textbf{$\alpha_A^* \to 0$ (attentional collapse):}
Occurs when $\zeta_\alpha R_A^{\text{surface}} \gg \gamma C_A$.
Then $R_0 \to 0$ and only $\mathcal{E}_0$ exists. Schema formation
is impossible regardless of training signal.

\textbf{$\delta_A^* = 0$ boundary:}
If the learning drive $f_{\text{learn}} T_A \eta_{\max}$ is
insufficient to overcome decay at $\sigma_A^*$, the only depth
fixed point is $\delta_A = 0$. This represents a training collapse
regime where the curriculum fails to induce any depth accumulation.
\end{boundary-case}

% --- Circular Dependency Check ---
\begin{circularity-check}
The proof uses the ODE definitions from Section~3.2 and
Proposition~4.1's definition of $\sigma_A^*$ (the transcritical
bifurcation result). This creates an apparent circularity:
Proposition~3.4 (equilibria) uses Proposition~4.1 ($\sigma_A^*$)
which depends on the equilibrium analysis.

\textbf{Resolution:} The dependency is unidirectional. Proposition~4.1
derives the existence of the transcritical bifurcation by analysing the
$\sigma_A$ ODE at equilibrium ($\dot{\sigma}_A = 0$) independently of
$\delta_A$ and $\alpha_A$. Proposition~3.4 uses that result to
characterise the full $3\times 3$ system after the bifurcation point
is established. Specifically:
\begin{itemize}[nosep]
\item The $\sigma_A$ equilibrium structure ($\sigma_A^* = (1 - 1/R_0)/\gamma_\sigma$)
  follows from setting $\dot{\sigma}_A = 0$ alone --- a one-dimensional
  analysis that does not depend on the $3\times 3$ Jacobian.
\item Proposition~3.4 then embeds this $\sigma_A^*$ into the full
  coupled system for the stability analysis via the $3\times 3$
  Jacobian.
\end{itemize}
This is a standard pattern in dynamical systems: first establish the
bifurcation structure of the slow variable ($\sigma_A$), then analyze
the full system stability around the resulting equilibria. No logical
circularity exists; the dependence is sequential and well-ordered.

\textbf{Forward dependencies:} Proposition~4.1 (bifurcation) references
the equilibrium existence but not the $3\times 3$ stability analysis.
Proposition~3.5 (numerical stability) uses the forward-invariant set
from Proposition~3.2 independently.
\end{circularity-check}
